This chapter covers the machine learning pipeline: data preparation, model
architectures, the diagnostic process that identified a fundamental feature
engineering flaw, and the final model performance after the pipeline rebuild.

% ═══════════════════════════════════════════════════════════
\section{Data Preparation}
% ═══════════════════════════════════════════════════════════

\subsection{Temporal Split}

A chronological split preserves the temporal ordering that the LSTM relies on.
A random split would leak future events into the training set, producing
artificially inflated validation metrics. The 180-day simulation period is
divided as follows:

\begin{table}[H]
\centering
\caption{Data split strategy}
\label{tab:data-split}
\begin{tabular}{@{}llrr@{}}
\toprule
\textbf{Split} & \textbf{Day Range} & \textbf{Events} & \textbf{Anomalies} \\
\midrule
Warm-up (discarded) & 1--21   & $\sim$28,000 & --- \\
Train (normal only)  & 22--126  & 142,254 & 0 \\
Validation           & 127--153 & 36,548  & 312 \\
Test                 & 154--180 & 36,059  & 175 \\
\bottomrule
\end{tabular}
\end{table}

The first 21 days are discarded as a warm-up period. During this window,
client buffers are still populating and velocity features
(\texttt{tx\_count\_24h}, \texttt{cumulative\_amount\_24h}) are unreliable.
The Autoencoder training set contains \textbf{normal events only} ---
anomalous events during the training window are excluded. This is the
standard approach for reconstruction-based anomaly detection: the model
learns what ``normal'' looks like, and any event it cannot reconstruct well
is flagged as anomalous.

A \texttt{StandardScaler} is fitted on the normal training data only and
applied to all splits. The scaler is saved alongside the model for use
during inference.

\subsection{Class Imbalance}

With an anomaly rate of 0.623\%, standard accuracy metrics are misleading
(a model that predicts ``normal'' for every event achieves 99.4\% accuracy).
AUC-ROC was chosen as the primary metric because it evaluates model
performance across all possible thresholds, independent of the class
distribution. Optuna hyperparameter searches use AUC-ROC as the objective
function.

% ═══════════════════════════════════════════════════════════
\section{Autoencoder}
% ═══════════════════════════════════════════════════════════

\subsection{Architecture}

The Autoencoder is a symmetric feedforward network that compresses the input
feature vector into a low-dimensional bottleneck and reconstructs it:

\begin{equation}
30 \xrightarrow{\text{encoder}} 20 \xrightarrow{} 10
\xrightarrow{\text{decoder}} 20 \xrightarrow{} 30
\end{equation}

The bottleneck dimension of 10 (one-third of input) provides enough capacity
to represent the structure of normal behavior without allowing the network to
learn an identity function. ReLU activations are used between layers. The
model is trained with Adam optimizer (learning rate $10^{-3}$), batch size
256, for 50 epochs with the best model selected by validation AUC.

\subsection{Anomaly Scoring}

The anomaly score is the \textbf{mean squared reconstruction error} between
the input vector $\mathbf{x}$ and its reconstruction $\hat{\mathbf{x}}$:

\begin{equation}
s_{\text{ae}} = \frac{1}{d} \sum_{i=1}^{d} (x_i - \hat{x}_i)^2
\end{equation}

Normal events, which the model has learned to reconstruct accurately, produce
low scores. Anomalous events, which deviate from the learned normal
distribution, produce high reconstruction errors.

% ═══════════════════════════════════════════════════════════
\section{LSTM}
% ═══════════════════════════════════════════════════════════

\subsection{Architecture}

The LSTM treats each client's transaction history as a temporal sequence.
It takes the last 10 enriched events as input and predicts what the next
event should look like:

\begin{equation}
\text{LSTM}(\text{input}=30, \text{hidden}=64, \text{layers}=2)
\rightarrow \text{Linear}(64, 30)
\end{equation}

The input at each time step is a 30-dimensional feature vector (the enriched
event). Two stacked LSTM layers process the sequence, and the final hidden
state feeds into a linear projection that predicts the 30 features of the
next event. The sequence length of 10 was chosen as a balance between
capturing sufficient behavioral context and keeping inference latency low.

\subsection{Anomaly Scoring}

The anomaly score is the mean squared prediction error between the LSTM's
predicted next event $\hat{\mathbf{y}}$ and the actual observed event
$\mathbf{y}$:

\begin{equation}
s_{\text{lstm}} = \frac{1}{d} \sum_{i=1}^{d} (y_i - \hat{y}_i)^2
\end{equation}

An event that matches the LSTM's prediction (consistent with the client's
behavioral history) produces a low score. An event that deviates from the
predicted pattern produces a high score. This approach is particularly
effective for detecting sequential anomalies --- patterns that appear normal
in isolation but are unusual given the client's recent activity.

\subsection{Sequence Construction}

Sequences are constructed per client using a sliding window: for each event
at position $t$ in a client's chronologically ordered history, the input
is events $[t{-}10, \ldots, t{-}1]$ and the target is event $t$. Clients
with fewer than 11 events are excluded from LSTM training. In the streaming
pipeline, the sequence buffer is maintained in Redis and updated after each
event.

% ═══════════════════════════════════════════════════════════
\section{Initial Training Failure and Root Cause Analysis}
\label{sec:pipeline-rebuild}
% ═══════════════════════════════════════════════════════════

The initial training attempt used the 74-feature enriched event schema.
Three Autoencoder configurations were evaluated via Optuna (71 features
with mean MSE, 34 behavioral features, 71 features with max/top-k
scoring). The best achieved AUC~0.577. The LSTM with multi-task prediction
heads achieved AUC~0.564. Both results were near random, triggering a
root cause investigation.

\subsection{Diagnosis}

Two independent root causes were identified:

\begin{enumerate}[leftmargin=2cm]
    \item \textbf{Undetectable scenarios.} Several of the original 13
          anomaly scenarios had stub implementations that did not produce
          measurable changes in the enriched feature space. For instance,
          \texttt{timing\_anomaly} shifted the event hour, but the enriched
          event contained no per-event hour feature --- only a profile-level
          hour distribution that absorbs single-event changes without
          producing a detectable signal.

    \item \textbf{Profile dominance.} The 74-feature enriched event was
          dominated by raw profile field dumps --- weekly snapshot means,
          standard deviations, and full distribution vectors that were
          identical across every event from the same client within the
          same week. The per-event signal was drowned out by profile noise.
          Anomalous events looked nearly identical to normal events from
          the same client because both carried the same profile values.
\end{enumerate}

\subsection{Resolution}

Three components were redesigned simultaneously:

\begin{itemize}[leftmargin=2cm]
    \item \textbf{Generator:} anomaly scenarios reduced from 13 to 9, with
          each remaining scenario validated for feature-space detectability
          (see Section~\ref{sec:scenarios}).
    \item \textbf{Enrichment:} complete redesign from profile dumps to
          contrast features. Feature count dropped from 74 to 30, with
          every feature measuring the deviation of the current event against
          the client's baseline (see Section~\ref{sec:enriched-event-design}).
    \item \textbf{Training data:} oracle profile shortcut introduced,
          providing exact client baselines without weekly granularity
          artifacts.
\end{itemize}

This was the single most impactful change in the project. The improvement
was entirely structural --- the same model architectures and hyperparameters
that produced AUC~0.57 on the 74-feature schema produced AUC~0.98+ on the
30-feature contrast schema.

% ═══════════════════════════════════════════════════════════
\section{Post-Rebuild Results}
% ═══════════════════════════════════════════════════════════

After the pipeline rebuild, both models were retrained on the redesigned
30-feature schema with branch-linked employees:

\begin{table}[H]
\centering
\caption{Model performance comparison}
\label{tab:model-performance}
\begin{tabular}{@{}lccc@{}}
\toprule
\textbf{Model} & \textbf{Initial AUC (74 feat.)} &
\textbf{Rebuilt AUC (30 feat.)} & \textbf{Final AUC (aligned)} \\
\midrule
Autoencoder & 0.577 & 0.9853 & 0.9714 \\
LSTM        & 0.564 & 0.9866 & 0.9355 \\
\bottomrule
\end{tabular}
\end{table}

The ``rebuilt'' column reflects training on the oracle profile shortcut (used
during development). The ``final'' column reflects retraining on the
production-aligned 30-feature schema where both streaming and batch
pipelines call the same \texttt{\_compute\_features()} function. The AUC
drop from rebuilt to final was expected: the Autoencoder's reconstruction
error is less sensitive to input source differences, while the LSTM lost
more because sequence modeling benefits from richer per-step representations.

\subsection{Per-Scenario Analysis}

\begin{table}[H]
\centering
\caption{Per-scenario detection performance (rebuilt models)}
\label{tab:per-scenario}
\begin{tabular}{@{}llcr@{}}
\toprule
\textbf{Scenario} & \textbf{Tier} & \textbf{AE AUC} & \textbf{LSTM AUC} \\
\midrule
Amount spike            & Easy   & 0.9998 & 0.9999 \\
Round amount            & Easy   & 1.0000 & 1.0000 \\
Duplicate virement      & Easy   & 0.9980 & --- \\
Frequency burst         & Medium & 0.9983 & 0.9125 \\
Cumulative threshold    & Medium & 0.9742 & 0.8540 \\
Timing anomaly          & Medium & 0.9710 & --- \\
Smurfing                & Hard   & 1.0000 & 0.9981 \\
Cheque structuring      & Hard   & 1.0000 & 0.9999 \\
Repeated client-employee & Hard  & 0.9983 & 0.9976 \\
\bottomrule
\end{tabular}
\end{table}

An unexpected finding: hard-tier scenarios produced higher AUC than
medium-tier scenarios. This inversion occurs because hard scenarios (smurfing,
cheque structuring) modify high-signal features like \texttt{z\_amount} and
\texttt{is\_near\_threshold} with large deviations, while medium scenarios
(frequency burst, cumulative threshold) produce subtler, multi-event patterns
that require more context to detect. Difficulty classification reflects
real-world detection challenge (deliberately subtle patterns), not model
sensitivity.

\subsection{Complementary Detection Profiles}

The two models exhibit complementary strengths that justify the dual-model
architecture:

\begin{itemize}[leftmargin=2cm]
    \item The \textbf{Autoencoder} excels at point anomalies visible in a
          single event (amount spike, round amount, threshold proximity).
          It produces high scores from the very first event, with no
          sequence history required.
    \item The \textbf{LSTM} excels at sequential patterns
          (repeated client-employee interactions, behavioral sequences).
          Its performance on medium-tier scenarios is lower because these
          patterns span multiple events and require sufficient context.
\end{itemize}

% ═══════════════════════════════════════════════════════════
\section{Framing for Production Deployment}
% ═══════════════════════════════════════════════════════════

These AUC values validate the system's ability to learn behavioral patterns
and serve predictions end-to-end. They are \textbf{not} a production
performance benchmark --- the models were trained on synthetic data whose
statistical properties approximate but do not replicate real Amen Bank
transaction distributions.

The engineering argument is more important than the absolute numbers: the
system is designed so that replacing synthetic data with real transactions
requires only re-running the training pipeline. The same
\texttt{\_compute\_features()} function, the same score fusion mechanism,
the same decision rules, and the same retraining infrastructure handle real
data without code changes. The pipeline is the deliverable, not the model
weights.

\subsection{Multi-Mechanism Convergence: Smurfing as Case Study}

Smurfing illustrates how the three detection mechanisms converge on a single
anomaly pattern. A smurfing scenario spans 21 days, but neither the
Autoencoder (which sees one event at a time) nor the LSTM (which sees a
10-event window) needs to observe the entire window to detect it.

Each smurfing event arrives at the Scoring Service already enriched with
velocity features computed from the rolling Redis buffer:
\texttt{near\_threshold\_count\_7d} accumulates across the smurfing window,
\texttt{cumulative\_amount\_24h} reflects recent activity, and
\texttt{z\_amount} captures the deviation from the client's baseline. These
features carry the behavioral \textit{history} into each event's feature
vector, regardless of the model's own window size.

The three mechanisms contribute distinct signals:

\begin{enumerate}[leftmargin=2cm]
    \item \textbf{Autoencoder:} each individual smurfing event is anomalous
          in isolation --- a salaried client whose retrait average is 300~DT
          suddenly performing 8,500~DT versements produces a high
          reconstruction error from \texttt{z\_amount},
          \texttt{amount\_to\_mean\_ratio}, and
          \texttt{is\_near\_threshold} alone.

    \item \textbf{LSTM:} the 10-event sequence shows escalating velocity
          features (\texttt{near\_threshold\_count\_7d} climbing from 0 to
          4, \texttt{tx\_count\_7d} rising) that deviate from the predicted
          normal trajectory. The LSTM does not need to see the full 21-day
          window --- it detects the \textit{acceleration} within its own
          context.

    \item \textbf{Decision Service Rule~2:} fires deterministically when
          the amount falls in the 8,000--9,999~DT band \textit{and} the
          client has $\geq$2 near-threshold transactions in 7 days.
          This rule requires no ML inference and catches structuring
          regardless of model scores.
\end{enumerate}

This convergence is not accidental --- it is a direct consequence of the
contrast feature design. Because each enriched event carries both its own
signal and the accumulated behavioral context from the rolling buffer, even
a stateless model (the Autoencoder) can participate in detecting a
multi-day pattern. The LSTM adds temporal sensitivity, and the regulatory
rules provide a deterministic safety net.